\phantomsection
\section*{Глава 1. Анализ предметной области}
\addcontentsline{toc}{section}{Глава 1. Анализ предметной области}
\stepcounter{section}

\subsection{Анализ предметной области сборки и продажи ПК}

Предметная область курсового проекта связана с деятельностью по закупке, хранению, учёту, сборке и продаже персональных компьютеров. Такая деятельность объединяет несколько взаимосвязанных процессов: работу со складом комплектующих, подбор совместимых компонентов, формирование готовых конфигураций, оформление заказов, взаимодействие с клиентами и получение аналитической информации по продажам.

Основой процесса является \textbf{складской учёт}. На складе могут находиться процессоры, материнские платы, видеокарты, оперативная память, накопители, блоки питания, корпуса, системы охлаждения и иные комплектующие. Для каждого товара необходимо хранить его наименование, категорию, количество, закупочную стоимость, продажную цену, характеристики и текущий статус наличия. При отсутствии автоматизированной системы ведение такого учёта вручную приводит к ошибкам, дублированию данных и трудностям при обновлении информации.

Следующий важный процесс --- \textbf{сборка персонального компьютера}. В отличие от обычной продажи единичных товаров, сборка ПК требует учёта совместимости комплектующих, контроля остатков и фиксации состава каждой конфигурации. При формировании одной сборки используются сразу несколько единиц товаров из разных категорий. После подтверждения сборки количество соответствующих комплектующих на складе должно автоматически уменьшаться, что требует поддержки связи между складским учётом и сущностью сборки.

Ещё одним ключевым направлением является \textbf{работа с заказами}. Клиент может выбрать готовую конфигурацию, заказать сборку по своим параметрам либо приобрести отдельные комплектующие. Для каждого заказа требуется сохранять сведения о клиенте, составе заказа, дате оформления, сумме, статусе исполнения и способе оплаты. Кроме того, важно обеспечить возможность отслеживания истории заказов и формирования отчётных данных по продажам.

В рамках данной предметной области можно выделить несколько \textbf{категорий пользователей}. Администратор отвечает за полное управление системой, настройку справочников и управление доступом. Менеджер работает с товарами, заказами и клиентами, а также может создавать и редактировать сборки. Клиент взаимодействует с приложением как пользователь, который просматривает ассортимент, выбирает товары и оформляет заказ. Следовательно, в системе должна быть реализована ролевая модель доступа.

Особенностью данной предметной области является необходимость обработки \textbf{аналитической информации}: пользователю системы важно не только добавлять и изменять записи, но и получать агрегированные данные --- общую сумму продаж за период, количество проданных товаров, наиболее востребованные комплектующие, остатки на складе, прибыль по заказам и по готовым сборкам.

Таким образом, предметная область включает совокупность процессов учёта товаров, формирования сборок, сопровождения заказов и анализа результатов продаж. Это делает разработку специализированного веб-приложения целесообразной и практически значимой.

\subsection{Обзор существующих веб-приложений и интернет-магазинов}

На современном рынке существует большое количество интернет-магазинов компьютерной техники и платформ для продажи комплектующих. К числу наиболее известных относятся крупные магазины, предлагающие широкий каталог товаров, фильтрацию по характеристикам, оформление заказов и базовые элементы личного кабинета пользователя. Такие решения ориентированы прежде всего на массовую торговлю и онлайн-продажи.

Несмотря на широкую распространённость подобных систем, они не всегда подходят для задач малого бизнеса или индивидуальной сборки ПК под заказ. Во многих случаях веб-сервисы ориентированы только на витрину товаров и не предоставляют полноценного механизма внутреннего учёта комплектующих, отслеживания фактических остатков или ведения истории использования компонентов в конкретных сборках.

Отдельно следует отметить \textbf{конфигураторы ПК}, позволяющие подобрать комплектующие по категориям. Их преимущество заключается в помощи пользователю при формировании конфигурации, однако подобные системы часто не интегрированы с внутренним складским учётом и механизмом управления заказами. В результате владелец бизнеса вынужден использовать сразу несколько разрозненных инструментов: отдельно таблицы для склада, отдельно мессенджеры или CRM для клиентов и отдельно сайт для публикации товаров.

Также существуют \textbf{универсальные CRM- и ERP-системы}, которые могут использоваться для учёта товаров и заказов. Их достоинством является наличие модулей для управления клиентами и продажами. Однако такие решения часто являются избыточными по функциональности, сложны в настройке и адаптации под конкретный сценарий сборки персональных компьютеров, а их использование может быть экономически невыгодным для небольших проектов.

Для наглядного сопоставления существующих программных продуктов проведён сравнительный анализ четырёх характерных представителей рынка: крупного интернет-магазина электроники DNS, специализированной сети «Регард», магазина-конфигуратора «НИКС» и универсальной торговой площадки «Ситилинк». Сравнение по существенным для предметной области критериям приведено в таблице~\ref{tab:competitors}.

\begin{table}[ht!]
\caption{Сравнительный анализ существующих веб-приложений}
\label{tab:competitors}
\footnotesize
\begin{tabularx}{\textwidth}{|>{\RaggedRight\arraybackslash}p{0.30\textwidth}|>{\centering\arraybackslash}X|>{\centering\arraybackslash}X|>{\centering\arraybackslash}X|>{\centering\arraybackslash}X|>{\centering\arraybackslash}X|}
\hline
\textbf{Критерий} & \textbf{DNS} & \textbf{Регард} & \textbf{НИКС} & \textbf{Сити-линк} & \textbf{Разра-баты-ваемое} \\ \hline
Каталог комплектующих & + & + & + & + & + \\ \hline
Конфигуратор сборки ПК & + & --- & + & --- & + \\ \hline
Контроль совместимости комплектующих & частично & --- & + & --- & + \\ \hline
Внутренний учёт остатков на складе & --- & --- & --- & --- & + \\ \hline
Ведение заказов в едином интерфейсе & + & + & + & + & + \\ \hline
Ролевая модель (админ/менеджер/клиент) & --- & --- & --- & --- & + \\ \hline
Аналитика по продажам и прибыли & --- & --- & --- & --- & + \\ \hline
Импорт/экспорт каталога (CSV, Excel) & --- & --- & --- & --- & + \\ \hline
Адаптация под малый бизнес & низкая & низкая & сред-няя & низкая & полная \\ \hline
Стоимость использования для бизнеса & нет & нет & нет & нет & собств. \\ \hline
\end{tabularx}
\end{table}

Результаты сравнения подтверждают сделанный ранее вывод: ни одно из рассмотренных решений не покрывает в комплексе задачи внутреннего учёта комплектующих, конфигурирования сборок, ведения заказов и аналитики. Существующие веб-приложения ориентированы прежде всего на витрину товаров для конечного покупателя и не предоставляют инструментов работы для администратора и менеджера небольшой компании по сборке ПК. Это окончательно обосновывает целесообразность разработки собственного специализированного веб-приложения.

\subsection{Анализ программных инструментов разработки}

Для реализации веб-приложения необходимо выбрать стек технологий, который обеспечит сбалансированное соотношение скорости разработки, надёжности эксплуатации и удобства последующего сопровождения. В таблице~\ref{tab:tech} приведено сравнение основных программных инструментов, рассматривавшихся в качестве кандидатов для серверной части, клиентской части и хранилища данных.

\begin{table}[ht!]
\caption{Сравнение программных инструментов разработки}
\label{tab:tech}
\small
\begin{tabularx}{\textwidth}{|>{\RaggedRight\arraybackslash}p{0.18\textwidth}|>{\RaggedRight\arraybackslash}p{0.20\textwidth}|>{\RaggedRight\arraybackslash}X|>{\RaggedRight\arraybackslash}X|}
\hline
\textbf{Инструмент} & \textbf{Назначение} & \textbf{Преимущества} & \textbf{Недостатки} \\ \hline
Django (Python) & Серверный фреймворк & Встроенная админ-панель, ORM, миграции, аутентификация, формы и валидация & Меньшая гибкость по сравнению с микрофреймворками \\ \hline
Flask (Python) & Серверный микрофреймворк & Минимализм, гибкость, простой запуск & Отсутствие готовой админки, много ручной работы \\ \hline
Node.js + Express & Серверная платформа & Высокая производительность, единый язык с фронтендом & Отсутствие готовой ORM, ручная сборка инфраструктуры \\ \hline
Laravel (PHP) & Серверный фреймворк & Зрелая экосистема, богатый набор пакетов & Менее распространён в российских учебных проектах \\ \hline
Bootstrap 5 & CSS-фреймворк & Готовая адаптивная сетка, набор компонентов, быстрый старт & Типовой вид интерфейса без кастомизации \\ \hline
Vanilla JS & Клиентская логика & Отсутствие зависимостей, простота отладки & Ручная реализация повторяющихся компонентов \\ \hline
React / Vue.js & Клиентские фреймворки & Компонентный подход, удобство для SPA & Усложняют сборку и SEO в учебном проекте \\ \hline
PostgreSQL & СУБД & Сложные запросы, надёжная транзакционная модель, ограничения целостности & Требует более тщательной настройки \\ \hline
SQLite & СУБД & Не требует сервера, удобна для разработки & Ограниченная пропускная способность \\ \hline
Docker & Контейнеризация & Унификация окружения, простота деплоя & Дополнительный слой инфраструктуры \\ \hline
Git + GitHub & Контроль версий & История изменений, ревью кода, распределённая работа & Требует дисциплины ведения коммитов \\ \hline
\end{tabularx}
\end{table}

На основании проведённого сравнения для реализации курсового проекта выбран следующий технологический стек:
\begin{itemize}
    \item \textbf{серверная часть} --- язык программирования Python~3.12 и фреймворк Django~5. Выбор обусловлен наличием встроенной административной панели, ORM для описания модели данных, готовых механизмов аутентификации, авторизации и валидации форм, а также большим объёмом учебной и справочной литературы~\cite{hawkins2022django};
    \item \textbf{клиентская часть} --- HTML5, CSS3, CSS-фреймворк Bootstrap~5 и Vanilla JavaScript. Такой набор позволяет получить адаптивный интерфейс, не усложняя структуру проекта дополнительными сборщиками и SPA-фреймворками;
    \item \textbf{хранилище данных} --- реляционная СУБД PostgreSQL~16 в качестве основной СУБД и SQLite на этапе локальной разработки. PostgreSQL обеспечивает поддержку ограничений целостности, транзакций и сложных аналитических запросов;
    \item \textbf{вспомогательные библиотеки} --- Django REST framework для построения внутреннего API, библиотека openpyxl и стандартный модуль \texttt{csv} для импорта и экспорта данных, библиотека Pillow для работы с изображениями товаров;
    \item \textbf{инфраструктура и инструменты разработки} --- система контроля версий Git и хостинг репозитория GitHub, среда разработки PyCharm Community / VS~Code, контейнеризация средствами Docker и Docker Compose, обратный прокси-сервер Nginx и WSGI-сервер Gunicorn для последующего развёртывания приложения.
\end{itemize}

Выбранный стек удовлетворяет требованиям учебного проекта по дисциплине «Разработка веб-приложений», обеспечивает быстрый старт разработки за счёт готовых компонентов Django и Bootstrap и позволяет в дальнейшем масштабировать приложение без полной переработки исходного кода.

\subsection{Формулировка цели, задач и требований к системе}

На основе проведённого анализа предметной области, обзора существующих решений и анализа программных инструментов сформулированы требования к разрабатываемому веб-приложению.

\textbf{Цель работы} --- разработка веб-приложения для учёта комплектующих, сборки и продажи персональных компьютеров с разграничением прав доступа администраторов, менеджеров и клиентов.

В системе должны быть реализованы функции \textbf{регистрации, аутентификации и авторизации} пользователей. Это необходимо для разграничения доступа к возможностям приложения в зависимости от роли пользователя. Администратор должен иметь полный доступ ко всем разделам системы, менеджер --- доступ к управлению товарами, заказами и сборками, клиент --- доступ к просмотру каталога, оформлению заказов и просмотру своих данных.

Важным требованием является наличие \textbf{CRUD-интерфейса} для основных сущностей приложения. Пользователь с соответствующими правами должен иметь возможность создавать, просматривать, изменять и удалять записи о товарах, сборках, заказах и пользователях.

Разрабатываемое приложение должно содержать \textbf{не менее трёх основных сущностей}. В рамках данного проекта такими сущностями являются:
\begin{itemize}
    \item пользователи (с ролями);
    \item категории и комплектующие;
    \item сборки ПК и их состав;
    \item заказы и позиции заказов.
\end{itemize}

Кроме хранения и редактирования данных, система должна обеспечивать \textbf{обработку и агрегирование} информации: вычисление суммы продаж за период, определение количества товаров на складе, расчёт прибыли по заказам и отображение наиболее востребованных комплектующих.

Отдельным требованием является \textbf{поддержка загрузки и выгрузки файлов} в форматах CSV и Excel для импорта начального каталога, экспорта отчётных данных и резервного копирования.

К \textbf{нефункциональным требованиям} относятся удобство пользовательского интерфейса, надёжность хранения данных, безопасность (хранение паролей в виде хэшей, защита от CSRF, XSS, SQL-инъекций), расширяемость и достаточная производительность при работе с каталогом, заказами и сборками.

Таким образом, разрабатываемая система должна представлять собой веб-приложение, объединяющее в себе механизмы учёта товаров, формирования сборок ПК, оформления заказов, разграничения доступа, загрузки и выгрузки файлов, а также аналитической обработки данных.
